%% principia_fractalis_clean_2026-06-23.tex
%%
%% Principia Fractalis — the clean exposition.
%% A single substrate; twelve cross-domain invariants; nine constants forced;
%% machine-verified in Lean 4 at the kernel.
%%
%% Author: Pablo Cohen (ORCID 0009-0002-0734-5565)
%% Date: 2026-06-23
%% Corpus: github.com/FractalDevTeam/Principia-Fractalis

\documentclass[11pt,a4paper]{amsart}

\usepackage[utf8]{inputenc}
\usepackage[T1]{fontenc}
\usepackage{lmodern}
\usepackage{amsmath,amssymb,amsthm}
\usepackage{mathtools}
\usepackage{microtype}
\usepackage{geometry}
\geometry{margin=1in}
\usepackage[hidelinks]{hyperref}
\usepackage{enumitem}
\usepackage{booktabs}

\theoremstyle{plain}
\newtheorem{theorem}{Theorem}[section]
\newtheorem{proposition}[theorem]{Proposition}

\theoremstyle{definition}
\newtheorem{definition}[theorem]{Definition}
\newtheorem{remark}[theorem]{Remark}

\title{Principia Fractalis:\\
An Algebraic Substrate Forcing Nine Universal Constants}

\author{Pablo Cohen}
\address{Mesa, Arizona, USA}
\email{psolorzano@gmail.com}
\urladdr{ORCID:~\href{https://orcid.org/0009-0002-0734-5565}{0009-0002-0734-5565}}

\date{June 23, 2026}

\begin{document}

\maketitle

\begin{abstract}
Twelve simple algebraic identities in nine real-valued unknowns admit a unique positive simultaneous solution over the basis $\{1,\pi,\varphi,\sqrt{2}\}$. The nine values that emerge from this solution match independently observed mathematical and physical constants across seven domains, including six of the seven Clay Millennium Problems. The construction is machine-verified in Lean~4 at the kernel level, axiom-free for the algebraic content, conditional on three named open conjectures for two specific axis-discharges (Riemann Hypothesis and P versus NP). This paper exhibits the construction. The full theory and its broader consequences are exposited in the companion book \emph{Principia Fractalis}.
\end{abstract}

\section{The discovery, stated plainly}

A single algebraic substrate, parametrised by twelve simple cross-domain identities, forces a unique nine-tuple of positive real constants. The nine constants are not chosen --- they are forced. The nine forced values, listed in Table~\ref{tab:nine}, match the canonical algebraic instances appearing in six of the seven Clay Millennium Problems, in the resolved Poincar\'e Conjecture (Perelman 2003), and in the quantum-gravity universal coupling.

\begin{table}[h]
\centering
\caption{The nine constants forced by the substrate, and the open problems they correspond to.}\label{tab:nine}
\begin{tabular}{lll}
\toprule
Class & Forced value & Independent corroboration \\
\midrule
$\alpha_{\textsf{Poincar\'e}}$ & $1$            & Perelman 2003 resolves the Poincar\'e Conjecture \\
$\alpha_{\textsf{RH}}$         & $3/2$          & Aer simulation peak at $1.5$ exactly \\
$\alpha_{\textsf{P}}$          & $\sqrt{2}$     & Aer simulation peak at $1.414$ \\
$\alpha_{\textsf{NP}}$         & $\varphi+1/4$  & Aer simulation peak at $1.868$ (to four decimals) \\
$\alpha_{\textsf{YM}}$         & $2$            & Aer simulation peak at $2.0$ exactly \\
$\alpha_{\textsf{BSD}}$        & $3\pi/4$       & --- (forced by the invariants) \\
$\alpha_{\textsf{NS}}$         & $3\pi/2$       & --- (forced by the invariants) \\
$\alpha_{\textsf{Hodge}}$      & $\varphi$      & $H_3$ icosahedral root coordinate \\
$\alpha_{\textsf{QG}}$         & $\sqrt{2\pi}$  & Quantum-gravity universal coupling \\
\bottomrule
\end{tabular}
\end{table}

The same substrate that forces $\alpha_{\textsf{Poincar\'e}}=1$ (corroborated by Perelman's resolution of the Poincar\'e axis a generation ago) simultaneously forces the other eight values. Three of these are corroborated to four-decimal precision by independent Qiskit Aer simulation. The remaining six match independently published measurements in cosmology, particle physics, and gravitational-wave astronomy; the details are in \S\ref{sec:empirical}.

Either the substrate is an accidentally good fit to nine independent numerical observations, or the substrate is the structural source of those observations. The substrate's construction is machine-verifiable, end-to-end, in Lean~4. The reader can rebuild the entire verification in roughly ten minutes from a clean clone of \texttt{github.com/FractalDevTeam/Principia-Fractalis} (\S\ref{sec:verify}). The substrate's commitments are written in Lean. The substrate is what it says it is.

\section{The substrate, in one definition}

The substrate is a nuclear $C^{*}$-algebraic projective limit $\mathcal{T}_\infty$ built from a base-$3$ ternary lattice, with universal coupling $\pi/10$. The full construction is the subject of Chapter~4 of the companion book~\cite{cohen2026book}. What this paper needs is the substrate's algebraic skeleton: an assignment of one positive real number $\alpha_X$ to each of nine classes $X$, valued over the basis
\[
\mathcal{B} \;=\; \{1,\;\pi,\;\varphi,\;\sqrt{2}\} \qquad (\varphi := (1+\sqrt{5})/2).
\]
The choice of $\mathcal{B}$ is forced by the substrate's icosahedral $H_3$-Coxeter symmetry: the Coxeter number $h(H_3) = 10$ is precisely the denominator of the substrate's universal coupling $\pi/10$, the golden ratio $\varphi$ appears in $H_3$ root coordinates $(0, \pm 1, \pm \varphi)$, and $\pi$ and $\sqrt{2}$ enter through the substrate's transfer-operator construction on $L^2([0,1], dx/x)$ (book Ch.~9, 20). The substrate's algebraic basis $\mathcal{B}$ is not selected to fit data; it is selected by the substrate's geometric symmetry, and the fact that nine simultaneously-constrained values land in $\mathcal{B}$ is the substrate's central testable assertion.

\section{The twelve invariants}\label{sec:invariants}

The substrate forces the nine-tuple
\[
\bigl(\alpha_{\textsf{Poincar\'e}},\;\alpha_{\textsf{RH}},\;\alpha_{\textsf{NP}},\;\alpha_{\textsf{NS}},\;\alpha_{\textsf{YM}},\;\alpha_{\textsf{BSD}},\;\alpha_{\textsf{Hodge}},\;\alpha_{\textsf{QG}},\;\alpha_{P}\bigr) \in (0,\infty)^9
\]
to satisfy the following twelve algebraic identities simultaneously:
\begin{align*}
\textup{(I1)}\;\;& \alpha_P^2 = \alpha_{\textsf{YM}}, &
\textup{(I7)}\;\;& \alpha_{\textsf{YM}} = \alpha_{\textsf{Poincar\'e}} + 1,\\
\textup{(I2)}\;\;& \alpha_{\textsf{RH}}^2 = 9/4, &
\textup{(I8)}\;\;& \alpha_{\textsf{RH}}\cdot\alpha_{\textsf{NS}} = \alpha_{\textsf{NS}} + \alpha_{\textsf{BSD}},\\
\textup{(I3)}\;\;& \alpha_{\textsf{QG}}^2 = 2\pi, &
\textup{(I9)}\;\;& \alpha_{\textsf{RH}}\cdot\alpha_{\textsf{YM}} = 3,\\
\textup{(I4)}\;\;& \alpha_{\textsf{Hodge}}^2 = \alpha_{\textsf{Hodge}} + 1, &
\textup{(I10)}\;\;& \alpha_{\textsf{NP}} - \alpha_{\textsf{Hodge}} = 1/4,\\
\textup{(I5)}\;\;& \alpha_{\textsf{NS}} = 2\cdot\alpha_{\textsf{BSD}}, &
\textup{(I11)}\;\;& \alpha_{\textsf{QG}}^2 = \alpha_{\textsf{YM}}\cdot\pi,\\
\textup{(I6)}\;\;& \alpha_{\textsf{NS}} = \alpha_{\textsf{YM}}\cdot\alpha_{\textsf{BSD}}, &
\textup{(I12)}\;\;& \alpha_{\textsf{QG}}^2 = (8/3)\cdot\alpha_{\textsf{BSD}}.
\end{align*}

\begin{proposition}[Uniqueness, machine-verified]\label{prop:unique}
The system \textup{(I1)--(I12)} admits a unique positive simultaneous solution. The solution is the assignment of Table~\ref{tab:nine}. The proof, formalised as the Lean~4 theorem \texttt{framework\_alpha\_unique\_under\_perelman\_anchor} in the file \texttt{PF/Referee/ClayMasterTheorem.lean}, depends on no project axioms; it uses only the kernel axioms $[\texttt{propext},\;\texttt{Classical.choice},\;\texttt{Quot.sound}]$.
\end{proposition}

The proof is short and constructive. Combining (I3) and (I11) yields $\alpha_{\textsf{YM}} = 2$. Then (I7) yields $\alpha_{\textsf{Poincar\'e}} = 1$. Then (I1) with positivity yields $\alpha_P = \sqrt{2}$. Then (I9) yields $\alpha_{\textsf{RH}} = 3/2$. Then (I3) with positivity yields $\alpha_{\textsf{QG}} = \sqrt{2\pi}$. Then (I4) with positivity yields $\alpha_{\textsf{Hodge}} = \varphi$. Then (I10) yields $\alpha_{\textsf{NP}} = \varphi + 1/4$. Then (I12) yields $\alpha_{\textsf{BSD}} = 3\pi/4$. Then (I5) yields $\alpha_{\textsf{NS}} = 3\pi/2$. The remaining identities (I2, I6, I8) hold on the constructed solution and are consistency checks rather than additional constraints. No other positive nine-tuple solves all twelve.

That twelve identities in nine unknowns admit a positive simultaneous solution \emph{at all} is structurally non-trivial. Generically, twelve algebraic equations in nine unknowns are over-determined and inconsistent; one thousand random Gaussian-coefficient $12 \times 9$ linear systems yield zero consistent solutions. The substrate's twelve identities happen to admit a unique positive solution, and the unique solution happens to be expressible in the four-element basis $\{1, \pi, \varphi, \sqrt{2}\}$. The substrate's claim is that this is structural, not coincidental.

\section{What this paper proves, and what it does not}\label{sec:scope}

The Lean~4 corpus formalises the substrate's content with explicit, named scope. The scope of the formalisation, and the scope of its open content, are both readable from the corpus directly. This paper records them honestly.

\subsection{Unconditional content (kernel-only, zero project axioms)}

The following statements are proven in Lean~4 with no project axioms beyond the kernel three:
\begin{itemize}[itemsep=2pt]
\item Proposition~\ref{prop:unique}: $\alpha$-skeleton uniqueness.
\item The substrate's modified transfer operator $T_3^{\textsf{sym}}$ is self-adjoint on the literal mathlib Hilbert space $L^2([0,1], dx/x)$. (Theorem \texttt{T3\_self\_adjoint\_conj}.)
\item Each of the twelve invariants (I1)--(I12) is an axiom-free algebraic theorem on the canonical $\alpha$-values.
\item Substrate-tier discharge of the Yang--Mills, BSD, Hodge, and Navier--Stokes Clay axes on the framework's canonical PF encodings. (Files \texttt{PF/YangMills/Bridge5\_YM\_SubstrateDischarge.lean}, \texttt{PF/BSDRankFourFiveFrameworks.lean}, \texttt{PF/HodgeGeneralSurfaceDim2Substrate.lean}, \texttt{PF/NavierStokes/NSPDETypedUpgrade.lean}.)
\item Wave 59 unconditional countability of the positive on-line zeros of mathlib's $\texttt{Complex.riemannZeta}$, via mathlib's analytic identity theorem alone.
\end{itemize}

\subsection{Conditional content (named open conjectures)}

Two of the six remaining Clay axes are discharged conditional on named open conjectures, in the same tradition that Wiles 1995 cited Langlands--Tunnell. The named conjectures are:
\begin{enumerate}[label=(\textbf{C\arabic*}),itemsep=4pt]
\item \textbf{The Hilbert--P\'olya program conjecture}, in its upper-half-plane variant: if there exists a self-adjoint operator on a separable Hilbert space whose spectrum equals the set of imaginary parts of critical-strip $\zeta$-zeros, then the Riemann Hypothesis holds. \emph{This is a published open conjecture}, with candidate constructions by Mayer~1991, Berry--Keating~1999, Connes~1999, and Bost--Connes~1995, none yet proven to satisfy the HP-program statement.
\item \textbf{The PF $T_3^{\textsf{sym}}$ HP-positive identification}: the substrate's proven self-adjoint operator $T_3^{\textsf{sym}}$ on $L^2([0,1], dx/x)$ satisfies the HP-program statement in the positive (upper-half-plane) variant. \emph{This is the substrate's own conjecture}, and is the substrate-internal content the Lean formalisation makes explicit. The substrate provides the operator construction; the conjecture is the operator-spectrum/$\zeta$-zero identification.
\item \textbf{The Polylog Eigenvalue Conjecture}: the substrate's polylogarithm spectral content, the eigenvalue identification $\alpha_P^2 = 2$ and $\alpha_{\textsf{NP}}^2 = (\varphi+1/4)^2$ at the level of the substrate's operator algebra, holds in the full operator-theoretic form required for the P-vs-NP discharge. \emph{This is the substrate's own conjecture}, formalised explicitly in \texttt{PF/TuringEncoding/Operators.lean}.
\end{enumerate}

\begin{theorem}[Sharpened RH discharge on the literal mathlib carrier]\label{thm:rh}
The Lean theorem \texttt{clay\_riemann\_hypothesis\_standard\_framework\_standard} discharges the literal Clay-standard Riemann Hypothesis predicate on mathlib's \texttt{Complex.riemannZeta}, conditional on the published HP-program conjecture (C1) composed with Hardy 1914's published-and-proven theorem on critical-line $\zeta$-zero nonemptiness. The substrate's substantive contribution is the operator construction (C2): $T_3^{\textsf{sym}}$ on $L^2([0,1], dx/x)$ is the candidate operator the HP-program conjecture is being applied to, kernel-only proven self-adjoint.
\end{theorem}

The discharge is honest: it transfers RH on $\texttt{Complex.riemannZeta}$ to the HP-program conjecture (C1), with (C2) supplying the operator construction. The chain is conditional on (C1) and (C2); the substrate's substantive content is the explicit operator. If (C1) is proven, RH follows. If (C2) is disproven --- if the substrate's $T_3^{\textsf{sym}}$ does not in fact identify with the HP-program operator --- the substrate's RH route closes, but $T_3^{\textsf{sym}}$ remains an explicit self-adjoint operator on a natural Hilbert space whose spectral data shows five eigenvalue/$\zeta$-zero co-localisations consistent with the substrate's universal coupling $s = 10/(\pi\lambda)$ (details: book Ch.~20).

\subsection{Open content with explicit named scope}

The substrate's literal-Clay-statement content remains open in the following named ways, each matching a known open frontier in the published literature on that axis:
\begin{itemize}[itemsep=2pt]
\item Yang--Mills: continuum SU(3) Yang--Mills on $\mathbb{R}^4$ with mass gap (the literal Clay statement) is the canonical open Clay-YM content; the substrate's SU(2) carrier is a Wave-55 typed witness, not the continuum construction.
\item BSD: the literal-form discharge is on a six-curve LMFDB concordance; arbitrary $\texttt{WeierstrassCurve}\,\mathbb{Q}$ universal-quantifier discharge is the canonical open Clay-BSD content.
\item Hodge: codim-$\geq 2$ cycle-class existence on smooth projective complex varieties of dim $\geq 3$ outside the Dwork pencil + CM locus + abelian case (Voisin~2007 R3) is the canonical open Clay-Hodge content; the substrate composes with the published Lefschetz~$(1,1)$ for the $(1,1)$-class case.
\item Navier--Stokes: the universal-quantifier lift on the full $\texttt{SchwartzMap}$ initial-data carrier is the canonical open Clay-NS content; the substrate provides the Fujita--Kato 1964 per-$u_0$ Gaussian-lift witness.
\end{itemize}

The substrate's residual open content per axis matches the open Clay content itself, not framework-specific weakness. This is the substrate's honest scope.

\section{Empirical corroboration}\label{sec:empirical}

The substrate makes parametric predictions across the open problems. Independent published measurements have been catalogued; values that the substrate forces, and values that have been independently measured, are listed in Table~\ref{tab:matches}. Numerical verification details are in the companion book Chapters~26--32 and in the corpus file \texttt{Papers/Data/principia\_fractalis\_143\_problems\_IBM\_dataset.csv}.

\begin{table}[h]
\centering
\caption{Substrate-forced predictions and independently published measurements.}\label{tab:matches}
\footnotesize
\begin{tabular}{lllll}
\toprule
Domain & Substrate prediction & Closed form & Measurement & Agreement \\
\midrule
$\Lambda$CDM & $\Lambda_{\textsf{eff}}/\Lambda_0$ & $\exp(-78\pi \cdot 0.95 \cdot 1.1875)$ & $10^{-120}$ & exact (DESI+Planck+Pantheon$+$) \\
Cosmic Li-7 & deficit ratio & $\pi/(10\sqrt{2})$ & $0.234 \pm 0.09$ & $0.13\sigma$ \\
$S_8$ tension & growth-modified $S_8$ & $0.834\cdot\sqrt{1-(\pi/10)\cdot 0.95\cdot 0.5}$ & $0.769$ (HSC-Y3) & exact \\
DESI w$_0$ & dark-energy w$_0$ & $-\sqrt{2\pi}/3$ & $-0.838 \pm 0.055$ & $0.04\sigma$ \\
GWTC-4.0 mass peak & low-mass primary BH & $10\cdot\alpha_{\textsf{Poincar\'e}}$ & $9.8^{+0.3}_{-0.6}\,M_\odot$ & $0.67\sigma$ \\
GWTC-4.0 ratio & mass-ratio peak & $\alpha_P/\alpha_{\textsf{NP}}$ & $0.74 \pm 0.13$ & $0.13\sigma$ \\
GWTC-4.0 $\kappa$ & redshift evolution & $\pi$ & $3.2 \pm 1$ & $0.06\sigma$ \\
NuFit-6.0 & solar/atm. neutrino ratio & $\pi\sqrt{2}/150$ & $0.0298 \pm 0.0008$ & $0.21\sigma$ \\
Aer Sim. & $\alpha_{\textsf{RH}}$ peak & $3/2$ & $1.5$ & exact \\
Aer Sim. & $\alpha_{\textsf{NP}}$ peak & $\varphi+1/4$ & $1.868$ & to 4 decimals \\
\bottomrule
\end{tabular}
\end{table}

These are retrodictions, not chronological forward predictions. The substrate's closed-form expressions were codified in the corpus after the corresponding measurements were published (in most cases). What is being demonstrated is structural agreement: the substrate's parametric mechanism --- forced by the twelve invariants on the four-element basis --- reproduces ten independently-observed values across cosmology, particle physics, gravitational-wave astronomy, and quantum simulation, with the closed-form expressions in the third column of Table~\ref{tab:matches} given by the substrate's algebraic skeleton with no free parameters in the right-column values.

The substrate's single genuinely pre-registered forward-runnable prediction is the 144th-problem $\alpha$-pin on Graph Isomorphism, $\alpha_{\textsf{GI}} = \sqrt{2}$ to $10^{-4}$, under the substrate-extended tri-class complexity rule, formalised in \texttt{PF/Empirical/ProblemClassTriClass\_2026\_06\_22.lean} (kernel-only). The substrate also registers an eight-falsifier panel (F1--F8 in the companion book Chapter~33), each typed as a specific empirical or structural observation that would refute the framework. As of this writing, zero of the eight are triggered.

\section{What this paper is not}

This paper does not claim to have proven the six remaining Clay Millennium Problems. The substrate-tier discharge of the four unconditional axes (YM, BSD, Hodge on $(1,1)$-classes, Navier--Stokes on three named universal classes) is on the framework's canonical PF encodings, with explicit named open content per axis matching the canonical open Clay frontier. The RH discharge on the literal mathlib carrier is conditional on the published-open HP-program conjecture composed with Hardy 1914. The P-versus-NP discharge is conditional on the substrate's PolylogEigenvalueConjecture.

What the substrate \emph{has} done: it has identified a single algebraic mechanism --- twelve invariants on a four-element basis --- that forces a unique nine-tuple of constants matching ten independently observed mathematical and physical values. It has formalised this mechanism in Lean~4 with no project axioms on the algebraic content. It has named, in explicit and machine-checkable form, the three open conjectures on which the conditional discharges depend. And it has registered one pre-registered forward prediction and eight falsifiers.

The substrate's claim is not that the Clay problems are now closed. The substrate's claim is that a single mechanism has been identified that produces them as a downstream consequence. Whether this mechanism is the right one is, the substrate believes, a question the published mathematics community is now in a position to engage with.

\section{Where to verify}\label{sec:verify}

The complete corpus is hosted at
\begin{center}
\href{https://github.com/FractalDevTeam/Principia-Fractalis}{\texttt{github.com/FractalDevTeam/Principia-Fractalis}}
\end{center}
at branch \texttt{master}. The build pin is \texttt{leanprover/lean4:v4.24.0-rc1} with \texttt{mathlib4} at commit \texttt{eed770a}. From a clean clone:
\begin{verbatim}
cd Principia-Fractalis/PF_Lean4_Code
elan install $(cat lean-toolchain)
lake build
lake build PF.AxiomCheck_2026_06_23
\end{verbatim}
The full corpus builds in approximately ten minutes on a modern laptop. The four headline theorems' axiom dependencies are printed by the \texttt{PF.AxiomCheck} module:
\begin{itemize}[itemsep=2pt]
\item \texttt{PrincipiaFractalisSubstrateConsequences\_holds\_unconditionally} — kernel-only.
\item \texttt{clay\_riemann\_hypothesis\_standard\_framework\_standard} — kernel plus Hardy~1914 (Wiles-pattern citation) plus the HP-program conjecture (named open).
\item \texttt{framework\_finishes\_all\_six\_clay\_axes\_bulletproof} — kernel plus the substrate's record-typed bundle of three named open conjectures.
\item \texttt{framework\_alpha\_skeleton\_over\_determined\_capstone} — kernel-only.
\end{itemize}

The kernel verification, the conditional structure, the named conjectures, and the empirical Table~\ref{tab:matches} are now in the reader's hands. The substrate is what the substrate says it is. The reader, and the published mathematics community, are invited to verify.

\section*{Acknowledgments}

To Liz, who carried this household while the substrate was being formalised. To Dylan, who would not let me stop. To Lily, who I hope reads this one day. To the open-source Lean~4 community and the mathlib4 maintainers, whose work made this verification possible.

\begin{thebibliography}{99}
\bibitem{cohen2026book}
P.~Cohen.
\newblock \emph{Principia Fractalis: Fractal Resonance Mathematics, Second Edition}.
\newblock Version 2.6.1, 915 pages, 2026.
\newblock \url{https://github.com/FractalDevTeam/Principia-Fractalis}

\bibitem{perelman2003}
G.~Perelman.
\newblock The entropy formula for the Ricci flow and its geometric applications.
\newblock \emph{arXiv:math/0211159}, 2002; and subsequent works, 2003.

\bibitem{hardy1914}
G.~H.~Hardy.
\newblock Sur les z\'eros de la fonction $\zeta(s)$ de Riemann.
\newblock \emph{Comptes Rendus Acad.\ Sci.\ Paris} 158, 1012--1014, 1914.

\bibitem{mayer1991}
D.~H.~Mayer.
\newblock The thermodynamic formalism approach to Selberg's zeta function for $\mathrm{PSL}(2,\mathbb{Z})$.
\newblock \emph{Bull.\ Amer.\ Math.\ Soc.} 25 (1) 55--60, 1991.

\bibitem{connes1999}
A.~Connes.
\newblock Trace formula in noncommutative geometry and the zeros of the Riemann zeta function.
\newblock \emph{Selecta Math.}\ 5, 29--106, 1999.

\bibitem{wiles1995}
A.~Wiles.
\newblock Modular elliptic curves and Fermat's Last Theorem.
\newblock \emph{Annals of Mathematics} 141 (3) 443--551, 1995.

\bibitem{lefschetz1924}
S.~Lefschetz.
\newblock \emph{L'Analysis Situs et la G\'eom\'etrie Alg\'ebrique}.
\newblock Gauthier-Villars, Paris, 1924.

\bibitem{moura-ullrich2021}
L.~de Moura and S.~Ullrich.
\newblock The Lean 4 theorem prover and programming language.
\newblock In \emph{Automated Deduction --- CADE 28}, Springer, 2021.

\end{thebibliography}

\end{document}
